% !TEX root = ../../../main.tex
% 来源：中学数学实验教材·第三册（高中数学）
% 章节：三角函数的图象和性质 / 正弦型曲线
% 原始文件：7.tex，行 1635-1658
% 类型：tikz
% ---
\begin{tikzpicture}[scale=1.3, >=latex]
\draw[->] (-2,0)--(7,0)node[right]{$x$};
\draw[->] (0,-2.5)--(0,2.5)node[right]{$y$};
\draw [domain=-pi/3:2*pi, samples=1000, dashed] plot(\x, {1.5*sin(3*\x r - pi r)});
\draw [domain=pi/3:pi, samples=1000, very thick] plot(\x, {1.5*sin(3*\x r - pi r)});
\node at (-.2,-.2){$O$};

\foreach \x/\xtext in {-2/-\frac{\pi}{3},-1/-\frac{\pi}{6},1/\frac{\pi}{6},3/\frac{\pi}{2},5/\frac{5\pi}{6},7/\frac{7\pi}{6},9/\frac{3\pi}{2},11/\frac{11\pi}{6}}
{
    \draw (\x*pi/6,0)node[below]{$\xtext$}--(\x*pi/6,.1);
}

\foreach \x/\xtext in {2/\frac{\pi}{3},4/\frac{2\pi}{3},6/\pi,8/\frac{4\pi}{3},10/\frac{5\pi}{3},12/2\pi}
{
    \draw (\x*pi/6,0)--(\x*pi/6,.1)node[above]{$\xtext$};
}

\foreach \x in {-1.5,-1,1,1.5}
{
    \draw(0,\x)node[left]{$\x$}--(.1,\x);
}


\end{tikzpicture}
